\documentclass[10pt,letterpaper]{article}
\usepackage[utf8]{inputenc}
\usepackage[spanish]{babel}
\usepackage{geometry}
\usepackage{tabularx}
\usepackage{booktabs}
\usepackage{amssymb}
\usepackage{fancyhdr}
\usepackage{titlesec}

% Configuración de márgenes
\geometry{top=2cm, bottom=2cm, left=1.5cm, right=1.5cm}

% Definición de casillas de verificación
\newcommand{\casilla}{\text{\Large$\square$}}
\newcommand{\casillaX}{\text{\Large$\boxtimes$}} % Usar para marcar con X

% =========================================================================
% DATOS DEL INFORME (Modifica estos campos para rellenar el documento)
% =========================================================================
\newcommand{\NoReferencia}{IPH-2026-XXXX}
\newcommand{\NoFolio}{FOLIO-12345}
\newcommand{\NoExpediente}{EXP-98765/2026}
\newcommand{\FechaPuesta}{20/05/2026}
\newcommand{\HoraPuesta}{11:43}

% Datos del Primer Respondiente
\newcommand{\ApellPRes}{Mireles}
\newcommand{\ApellSRes}{Alcántara}
\newcommand{\NomRes}{Abraham Apolinar}
\newcommand{\AdscripcionRes}{FCPyS - UNAM}
\newcommand{\CargoRes}{Oficial de Control}

% =========================================================================

\pagestyle{fancy}
\fancyhf{}
\fancyhead[L]{\textbf{SISTEMA NACIONAL DE SEGURIDAD PÚBLICA}}
\fancyhead[R]{IPH-2019 | HECHO PROBABLEMENTE DELICTIVO}
\fancyfoot[C]{Página \thepage}

\titleformat{\section}{\normalfont\large\bfseries\fillwithlines{1em}}{}{0pt}{}[\titlerule]

\begin{document}

\begin{center}
    {\LARGE \textbf{INFORME POLICIAL HOMOLOGADO}} \\
    \vspace{0.2cm}
    {\large \textbf{SECCIÓN 1. PUESTA A DISPOSICIÓN}}
\end{center}

\vspace{0.4cm}

% Tabla de Control Superior
\begin{tabularx}{\textwidth}{|X|X|X|}
\hline
\textbf{NO. DE REFERENCIA:} & \textbf{NO. DE FOLIO (SISTEMA):} & \textbf{NO. EXPEDIENTE:} \\
\NoReferencia & \NoFolio & \NoExpediente \\
\hline
\end{tabularx}

\vspace{0.6cm}

\subsubsection*{Apartado 1.1 Fecha y hora de la puesta a disposición}
\begin{tabularx}{\textwidth}{|X|X|}
\hline
\textbf{Fecha (DDMMAAAA):} \FechaPuesta & \textbf{Hora (HH:MM):} \HoraPuesta \\
\hline
\end{tabularx}

\vspace{0.6cm}

\subsubsection*{Anexos Entregados (Señale con una "X" e indique la cantidad)}
\begin{tabularx}{\textwidth}{|l|X|l|l|X|}
\hline
\casilla & Anexo A. Detención(es) & & \casilla & Anexo E. Entrevistas \\ \hline
\casilla & Anexo B. Informe del uso de la fuerza & & \casilla & Anexo F. Entrega-recepción del lugar \\ \hline
\casilla & Anexo C. Inspección de vehículo & & \casilla & Anexo G. Continuación de la narrativa \\ \hline
\casilla & Anexo D. Inventario de armas y objetos & & \casilla & No se entregan Anexos \\ 
\hline
\end{tabularx}

\vspace{0.6cm}

% Sección de firmas de puesta a disposición
\begin{tabularx}{\textwidth}{|X|X|}
\hline
\multicolumn{2}{|c|}{\textbf{DATOS DE QUIEN REALIZA LA PUESTA A DISPOSICIÓN (PRIMER RESPONDIENTE)}} \\
\hline
\textbf{Primer Apellido:} \ApellPRes & \textbf{Segundo Apellido:} \ApellSRes \\
\hline
\textbf{Nombre(s):} \NomRes & \textbf{Adscripción:} \AdscripcionRes \\
\hline
\textbf{Cargo/Grado:} \CargoRes & \textbf{Firma:} \\[0.8cm]
\hline
\end{tabularx}

\vspace{0.8cm}
\pagebreak

% =========================================================================
% SECCIÓN 2. PRIMER RESPONDIENTE Y UBICACIÓN
% =========================================================================
\section*{SECCIÓN 2 Y 3. CONOCIMIENTO DEL HECHO Y UBICACIÓN}

\subsubsection*{Apartado 3.1 Conocimiento del hecho por el primer respondiente}
¿Cómo se enteró del hecho? \\
\begin{tabularx}{\textwidth}{XXXX}
\casilla~Denuncia & \casilla~Flagrancia & \casilla~Localización & \casilla~Descubrimiento \\
\casilla~Mandamiento judicial & \casilla~Llamada de emergencia & \casilla~Aportación & No. 911: \underline{\hspace{3cm}}
\end{tabularx}

\vspace{0.6cm}

\subsubsection*{Apartado 4.1 Ubicación geográfica del lugar de la intervención}
\begin{tabularx}{\textwidth}{|X|X|X|}
\hline
\multicolumn{3}{|l|}{\textbf{Calle / Tramo carretero:}} \\ \hline
\textbf{No. Exterior:} & \textbf{No. Interior:} & \textbf{Código Postal:} \\ \hline
\multicolumn{2}{|l|}{\textbf{Colonia / Localidad:}} & \textbf{Municipio / Demarcación:} \\ \hline
\multicolumn{3}{|l|}{\textbf{Entidad federativa:}} \\ \hline
\textbf{Latitud:} & \textbf{Longitud:} & \textbf{Referencias:} \\ \hline
\end{tabularx}

\vspace{0.8cm}

% =========================================================================
% SECCIÓN 5. NARRATIVA DE LOS HECHOS
% =========================================================================
\section*{SECCIÓN 5. NARRATIVA DE LOS HECHOS}
\textit{Relate cronológicamente las acciones realizadas durante su intervención desde el conocimiento del hecho hasta la puesta a disposición. Responda de manera descriptiva a las preguntas: ¿Quién?, ¿Qué?, ¿Cómo?, ¿Cuándo? y ¿Dónde?.}

\vspace{0.4cm}
\begin{tabularx}{\textwidth}{|X|}
\hline
\textbf{Descripción de los hechos y actuación de la autoridad:} \\[10cm]
\hline
\end{tabularx}

\vspace{0.8cm}
\pagebreak

% =========================================================================
% ANEXO A. DETENCIONES (EJEMPLO DE ESTRUCTURA DE ANEXO)
% =========================================================================
\section*{ANEXO A. DETENCIÓN(ES)}
\subsubsection*{Apartado A.2 Datos generales de la persona detenida}

\begin{tabularx}{\textwidth}{|X|X|X|}
\hline
\textbf{Primer Apellido:} & \textbf{Segundo Apellido:} & \textbf{Nombre(s):} \\ \hline
\textbf{Apodo o alias:} & \textbf{Nacionalidad:} \casilla~Mexicana ~ \casilla~Extranjera & \textbf{Sexo:} \casilla~F ~ \casilla~M \\ \hline
\textbf{Fecha de Nacimiento:} & \textbf{Edad:} & \textbf{Identificación:} \casilla~INE \casilla~Pasaporte \\ \hline
\end{tabularx}

\vspace{0.6cm}
\subsubsection*{Apartado A.4 Constancia de lectura de derechos}
\begin{enumerate}
    \item Usted tiene derecho a conocer el motivo de su detención.
    \item Usted tiene derecho a guardar silencio.
    \item Usted tiene derecho a declarar, asistido por un defensor ante la autoridad competente.
    \item Usted tiene derecho a un defensor público si no puede o no quiere designar uno privado.
    \item Usted tiene derecho a que un familiar o persona de confianza conozca los hechos de su detención.
\end{enumerate}

\vspace{0.5cm}
\begin{center}
\begin{tabular}{p{7cm}}
\hline
\centering \textbf{Firma o Huella de la persona detenida}
\end{tabular}
\end{center}

\end{document}
